% ============ 第八章 §4 不分歧扩张和完全分歧扩张 (原书 pp.294-300) ============
\section{不分歧扩张和完全分歧扩张}\label{sec:ch8-s4}

本节中设 $(F,P_F)\subseteq(E,P_E)$ 是完备离散赋值域的有限扩张. 以 $O_E,\mathfrak{p}_E$ 表示 $E$ 的赋值环和素理想, $O_F,\mathfrak{p}_F$ 表示 $F$ 的赋值环和素理想, 而 $\overline{E}=O_E/\mathfrak{p}_E$ 和 $\overline{F}=O_F/\mathfrak{p}_F$ 分别表示 $E$ 和 $F$ 的剩余类域. 令 $e=e(P_E/P_F)$, $f=f(P_E/P_F)$, 则 $ef=n=[E:F]$. 如果 $e=1$（从而 $f=n$）, 称 $E/F$ 为\emph{不分歧扩张}. 如果 $f=1$（从而 $e=n$）, 称 $E/F$ 为\emph{完全分歧扩张}. 本节要刻画这些扩张的结构.

\begin{theorem}\label{thm:8.17}
(1) 设 $E/F$ 是不分歧扩张. 如果 $\overline{E}=\overline{F}(\alpha_0)$, 取元素 $\alpha\in O_E$, 使得 $\bar{\alpha}=\alpha_0$, 则 $E=F(\alpha)$, 并且若 $f(x)$ 是 $\alpha$ 在 $F$ 上的极小多项式, 则 $\overline{f}(x)$ 是 $\bar{\alpha}$ 在 $\overline{F}$ 上的极小多项式.

(2) 若 $E=F(\alpha)$, $\alpha\in O_E$, $g(x)$ 是 $O_F[x]$ 中首 $1$ 多项式, $g(\alpha)=0$. 如果 $\overline{g}(x)$（在 $\overline{F}$ 的代数闭包 $\overline{\Omega}$ 中）没有重根, 则 $E/F$ 是不分歧扩张.
\end{theorem}

\begin{proof}
(1) 设 $f(x)$ 是 $\alpha$ 在 $F$ 上的极小多项式. 由于 $\alpha\in O_E$, 从而 $f(x)$ 是 $O_F[x]$ 中首 $1$ 多项式, 并且 $\deg\overline{f}(x)=\deg f$, $\overline{f}(\alpha_0)=\overline{f}(\bar{\alpha})=0$. 于是
\[
  \deg\overline{f}\geqslant[\overline{F}(\alpha_0):\overline{F}]=[\overline{E}:\overline{F}]=[E:F],
\]
\[
  \deg f=[F(\alpha):F]\leqslant[E:F].
\]
但是 $\deg\overline{f}=\deg f$. 这就表明上面的不等式均为等式, 即 $E=F(\alpha)$, $\deg\overline{f}=[\overline{F}(\bar{\alpha}):\overline{F}]$, 所以 $\overline{f}(x)$ 为 $\bar{\alpha}$ 在 $\overline{F}$ 上的极小多项式.

(2) 不妨设 $g(x)$ 是 $\alpha$ 在 $F$ 上的极小多项式, 则 $\overline{g}(\bar{\alpha})=0$. 由于 $g(x)$ 在 $F$ 上不可约, 因此 $\overline{g}(x)$ 必为 $\overline{F}[x]$ 中某个不可约多项式的方幂（Hensel 引理）. 但是 $\overline{g}(x)$ 无重根, 从而 $\overline{g}(x)$ 必为 $\overline{F}[x]$ 中不可约多项式. 于是
\begin{align*}
  [E:F]&=[F(\alpha):F]=\deg g(x)=\deg\overline{g}\\
  &=[\overline{F}(\bar{\alpha}):\overline{F}]\leqslant[\overline{E}:\overline{F}]\leqslant[E:F],
\end{align*}
这就表明 $[E:F]=[\overline{E}:\overline{F}]$, 即 $E/F$ 是不分歧扩张.
\end{proof}

\begin{lemma}\label{lem:ch8-3}
设 $F\subseteq E$ 是完备离散赋值域的有限扩张, 并且 $\overline{F}$ 是有限域.

(1) 对于 $F\leqslant M\leqslant E$, 则 $E/F$ 不分歧当且仅当 $E/M$ 和 $M/F$ 均不分歧.

(2) 若 $K/F$ 是有限扩张, $E/F$ 不分歧, 则 $KE/K$ 不分歧.

(3) 若 $E_1/F$, $E_2/F$ 均不分歧, 则 $E_1E_2/F$ 不分歧.
\end{lemma}

\begin{proof}
(1) 由定义易知 $e(E/F)=e(E/M)e(M/F)$. 由此即推出结论.

(2) 由 $E/F$ 不分歧可知 $E=F(\alpha)$, 其中 $\alpha\in O_E$, 并且对于 $\alpha$ 在 $F$ 上的极小多项式 $f(x)$, $\overline{f}(x)$ 是 $\overline{F}[x]$ 中不可约多项式（定理 \ref{thm:8.17} 的 (1)）. 由于 $\overline{F}$ 是有限域, 可知 $\overline{f}(x)$（在 $\overline{F}$ 的代数闭包中）没有重根. 于是 $KE=KF(\alpha)=K(\alpha)$, $f(x)\in O_F[x]\subseteq O_K[x]$, $\overline{f}(x)\in\overline{F}[x]\subseteq\overline{K}[x]$ 没有重根, 所以由定理 \ref{thm:8.17} 的 (2) 即知 $KE/K$ 不分歧.

(3) 由 (2) 知 $E_1E_2/E_1$ 不分歧, 再由 (1) 即知 $E_1E_2/F$ 不分歧.
\end{proof}

引理 \ref{lem:ch8-3} 可用于 $F=\Q_p$（$p$ 进数域）的情形, 因为 $\Q_p$ 的剩余类域 $\Z_p/p\Z_p\cong\Z/p\Z$ 是有限域 $\F_p$. 以下固定 $\Q_p$ 的一个代数闭包 $\Omega_p$. $\Q_p$ 的有限扩张均是指 $\Omega_p$ 的子域.

\begin{theorem}\label{thm:8.18}
(1) 设 $\alpha$ 是 $\Omega_p$ 中的 $l$ 次本原单位根, 并且 $l$ 与 $p$ 互素. 以 $n$ 表示满足 $p^n\equiv 1\ (\mathrm{mod}\ l)$ 的最小正整数, 则 $\Q_p(\alpha)$ 是 $\Q_p$ 的 $n$ 次不分歧扩张, 并且这是伽罗瓦扩张, 其伽罗瓦群 $\Gal(\Q_p(\alpha)/\Q_p)$ 是由自同构 $\sigma$ 生成的 $n$ 次循环群, 其中 $\sigma(\alpha)=\alpha^p$.

(2) 对每个 $n\geqslant 1$, $\Q_p$ 都存在惟一的 $n$ 次不分歧扩张.

(3) 设 $E/\Q_p$ 是有限扩张, 则存在中间域 $F$, 使得 $F/\Q_p$ 为不分歧, 而 $E/F$ 是完全分歧.
\end{theorem}

\begin{proof}
(1) 由于 $\alpha$ 是 $f(x)=x^l-1\in\Z_p[x]$ 的根, 并且由 $(l,p)=1$ 可知 $\overline{f}(x)$ 在 $\F_p[x]$ 中没有重根. 于是 $\Q_p(\alpha)/\Q_p$ 不分歧. 有限域 $\F_p(\bar{\alpha})/\F_p$ 的扩张次数等于使 $\bar{\alpha}^{p^n}=\bar{\alpha}$ 成立的最小正整数 $n$, 而由 $\overline{f}(x)$ 没有重根可知 $\bar{\alpha}$ 仍是 $l$ 次本原单位根. 因此 $n$ 即为满足 $p^n\equiv 1\ (\mathrm{mod}\ l)$ 的最小正整数. 于是 $[\Q_p(\alpha):\Q_p]=[\F_p(\bar{\alpha}):\F_p]=n$.

进而, $\alpha$ 的某个共轭元素仍是 $l$ 次本原单位根, 从而必为 $\alpha^m$, 其中 $(m,l)=1$. 于是 $\alpha$ 的共轭元素均属于 $\Q_p(\alpha)$, 即 $\Q_p(\alpha)/\Q_p$ 为伽罗瓦扩张. 设 $\tau$ 是伽罗瓦群中的元素, 则 $\tau(\alpha)=\alpha^m$, 其中 $(m,l)=1$. 由 $\tau$ 诱导出有限域 $n$ 次扩张 $\F_p(\bar{\alpha})/\F_p$ 的自同构 $\bar{\tau}$, 其中 $\bar{\tau}(\bar{\alpha})=\bar{\alpha}^m$. 但是 $\F_p(\bar{\alpha})/\F_p$ 的伽罗瓦群是由 $\bar{\alpha}\mapsto\bar{\alpha}^p$ 生成的 $n$ 阶循环群, 于是 $\bar{\alpha}^m=\bar{\alpha}^{p^t}\ (0\leqslant t\leqslant n-1)$. 但是 $\alpha^m$ 和 $\alpha^{p^t}$ 都是 $\Q_p(\alpha)$ 中的 $p^n-1$ 次单位根, 而所有 $p^n-1$ 次单位根加上 $0$ 应当是 $\Q_p(\alpha)$ 的赋值环对其素理想的陪集完全代表系. 所以由 $\bar{\alpha}^m=\bar{\alpha}^{p^t}$ 可知 $\alpha^m=\alpha^{p^t}$. 这就表明自同构 $\tau$ 只有 $n$ 个可能性: $\tau(\alpha)=\alpha^{p^t}\ (0\leqslant t\leqslant n-1)$. 但是 $\Q_p(\alpha)/\Q_p$ 应当有 $n=[\Q_p(\alpha):\Q_p]$ 个自同构, 因此 $\tau_t(\alpha)=\alpha^{p^t}\ (0\leqslant t\leqslant n-1)$ 都是 $\Q_p(\alpha)/\Q_p$ 的自同构. 从而整个伽罗瓦群是由 $\tau_1$（$\tau_1(\alpha)=\alpha^p$）生成的 $n$ 阶循环群.

(2) 对每个 $n\geqslant 1$, 取 $\alpha$ 为 $\Omega_p$ 中的 $p^n-1$ 次本原单位根（$x^{p^n-1}-1$ 在 $\Omega_p$ 中没有重根, 全部 $p^n-1$ 个根应当是乘法循环群, 其生成元就是 $p^n-1$ 次本原单位根）. 由 (1) 知 $\Q_p(\alpha)$ 即是 $\Q_p$ 的 $n$ 次不分歧扩张. 再证惟一性: 若 $E$ 和 $F$ 均是 $\Q_p$ 的 $n$ 次不分歧扩张, 则 $\overline{E}$ 和 $\overline{F}$ 均是 $\F_p$ 的 $n$ 次扩张, 于是 $\overline{E}=\overline{F}$. 由引理 \ref{lem:ch8-3} 知 $EF$ 也为 $\Q_p$ 的不分歧扩张. 于是
\[
  [EF:\Q_p]=[\overline{EF}:\F_p]=[\overline{E}\:\overline{F}:\F_p]=[\overline{E}:\F_p]=n.
\]
所以 $[EF:F]=[EF:\Q_p]/[F:\Q_p]=1$, 即 $EF=F$, 因此 $E\subseteq F$. 再由 $[E:\Q_p]=[F:\Q_p]$ 可知 $E=F$. 这就表明对每个 $n\geqslant 1$, $\Q_p$ 恰好存在惟一的 $n$ 次不分歧扩张.

(3) 以 $F$ 表示 $E$ 中所有对 $\Q_p$ 不分歧的子域的合成, 则 $F/\Q_p$ 不分歧. 易知 $E/F$ 完全分歧.
\end{proof}

\begin{example}
对 $\Q_p$ 的每个二次扩域 $K$, $ef=2=[K:\Q_p]$. 所以或者 $e=1$, $f=2$（不分歧）, 或者 $e=2$, $f=1$（完全分歧）.

当 $p\geqslant 3$ 时, $\Q_p$ 共有 $3$ 个二次扩域 $\Q_p(\sqrt{a})$, 其中 $a=c,cp$ 和 $p$, 而 $1\leqslant c\leqslant p-1$, $\left(\frac{c}{p}\right)=-1$. 由于 $x^2-c$ 在 $\F_p$ 中不可约, 可知 $\Q_p(\sqrt{c})$ 是 $\Q_p$ 的不分歧扩张. 由惟一性知 $\Q_p(\sqrt{p})$ 和 $\Q_p(\sqrt{cp})$ 是 $\Q_p$ 的完全分歧扩张.

当 $p=2$ 时, $\Q_2$ 共有 $7$ 个二次扩域 $\Q_2(\sqrt{a})$, 其中 $a=3,5,7,2,6,10,14$. $\Q_2$ 的二次不分歧扩张应当是添加 $2^2-1=3$ 次本原单位根, 即添加 $x^2+x+1$ 的根 $\omega=\frac{1}{2}(-1+\sqrt{-3})$. 所以 $\Q_2(\omega)=\Q_2(\sqrt{-3})=\Q_2(\sqrt{5})$ 是 $\Q_2$ 的不分歧扩张. 而 $\Q_2$ 的其余 $6$ 个二次扩域均是完全分歧的. 这些域的完全分歧性也可由下面定理得出.
\end{example}

现在刻画完全分歧扩张.

\begin{theorem}\label{thm:8.19}
设 $E/F$ 是完备离散赋值域的有限扩张, $\pi$ 是 $E$ 的一个素元.

(1) 若 $E/F$ 是完全分歧的, 则 $E=F(\pi)$, 并且 $\pi$ 在 $F$ 上的最小多项式为 Eisenstein 多项式.

(2) 反之, 若 $E=F(\alpha)$, 并且 $\alpha$ 在 $F$ 上的最小多项式是 Eisenstein 多项式, 则 $E/F$ 是完全分歧扩张, 并且 $\alpha$ 是 $E$ 的一个素元.
\end{theorem}

\begin{proof}
(1) $\pi$ 在 $F$ 上的最小多项式为 $O_F[x]$ 中多项式 $f(x)=x^m+a_1x^{m-1}+\cdots+a_m$. 以 $v_F$ 和 $v_E$ 分别表示 $F$ 和 $E$ 中的标准指数赋值, $v$ 是 $v_F$ 在 $F$ 的代数闭包 $\Omega$ 中的（惟一）扩充, 则 $1=v_E(\pi)=nv(\pi)$, 其中 $n=[E:F]$（因为 $E/F$ 完全分歧, 从而 $e(E/F)=n$, 而 $v_E=nv|_E$）. 设 $\pi_1,\cdots,\pi_m$ 是 $f(x)$ 的 $m$ 个根（$\pi_1=\pi$）, 则 $v(\pi_i)\ (1\leqslant i\leqslant m)$ 均相等（它们是 $F$-共轭元素）. 于是
\[
  v_F(a_m)=v(a_m)=v(\pi_1)+\cdots+v(\pi_m)=mv(\pi)=\frac{m}{n}.
\]
但是 $v_F(a_m)\in\Z$. 因此 $m\geqslant n$. 另一方面, $n=[E:F]\geqslant[F(\pi):F]=m$. 这表明 $n=m$ 并且 $E=F(\pi)$. 由于 $a_i$ 是 $\pi_1,\cdots,\pi_m$ 的初等对称函数, 而 $v(\pi_i)=v(\pi)>0$, 因此 $v_F(a_i)=v(a_i)>0$, 即 $v_F(a_i)\geqslant 1\ (1\leqslant i\leqslant m)$. 再由 $v_F(a_m)=\frac{m}{n}=1$, 即知 $f(x)$ 是 Eisenstein 多项式.

(2) 设 $E=F(\alpha)$, $\alpha$ 在 $F$ 上的最小多项式 $f(x)=x^n+a_1x^{n-1}+\cdots+a_n$ 是 Eisenstein 多项式, 即 $v_F(a_i)\geqslant 1\ (1\leqslant i\leqslant n)$, $v_F(a_n)=1$. $f(x)$ 的牛顿折线即是从 $(n,v_F(1))=(n,0)$ 到 $(0,v_F(a_n))=(0,1)$ 的连线, 可知 $v(\alpha)=\frac{1}{n}$. 从而
\[
  e(E/F)=[v(E^{\times}):v_F(F^{\times})]\geqslant\left[\frac{1}{n}\Z:\Z\right]=n.
\]
这就表明 $e(E/F)=n$, 即 $E/F$ 完全分歧. 而 $v_E(\alpha)=e(E/F)v(\alpha)=n\cdot\frac{1}{n}=1$, 从而 $\alpha$ 为 $E$ 的素元.
\end{proof}

\begin{example}
当 $p$ 为奇素数时, $x^2-cp$ 和 $x^2-p$ 均是 Eisenstein 多项式（$1\leqslant c\leqslant p-1$, $\left(\frac{c}{p}\right)=-1$）. 所以 $\Q_p(\sqrt{p})$ 和 $\Q_p(\sqrt{cp})$ 都是 $\Q_p$ 的完全分歧扩张. 对于 $p=2$, 同样可知 $\Q_2(\sqrt{a})\ (a=2,6,10,14)$ 是 $\Q_2$ 的完全分歧扩张. 进而对 $\Q_2(\sqrt{3})$ 和 $\Q_2(\sqrt{7})$, 我们有 $\Q_2(\sqrt{3})=\Q_2(1+\sqrt{3})$, $\Q_2(\sqrt{7})=\Q_2(1+\sqrt{7})$. 而 $1+\sqrt{3}$ 和 $1+\sqrt{7}$ 在 $\Q_2$ 上的极小多项式 $x^2-2x-2$ 和 $x^2-2x-6$ 都是对 $p=2$ 的 Eisenstein 多项式. 所以 $\Q_2(\sqrt{3})$ 和 $\Q_2(\sqrt{7})$ 也是 $\Q_2$ 的完全分歧扩张.
\end{example}

\begin{example}
（分圆扩张）设 $p$ 为素数, $n=p^l\cdot m$, $l\geqslant 1$, $(p,m)=1$. 令 $\alpha$ 为 $\Omega_p$ 中 $n$ 次本原单位根, $F=\Q_p(\alpha)$. 我们计算 $F/\Q_p$ 的分歧指数 $e$ 和剩余类域次数 $f$, 而 $[F:\Q_p]=ef$.

令 $M_1=\Q_p(\beta)$, 其中 $\beta=\alpha^{p^l}$ 是 $m$ 次本原单位根, 则 $M_1$ 是 $F/\Q_p$ 的中间域. 由 $(p,m)=1$ 可知 $M_1/\Q_p$ 是不分歧扩张, $[M_1:\Q_p]=s$, 其中 $s$ 是满足 $p^s\equiv 1\ (\mathrm{mod}\ m)$ 的最小正整数. 另一方面, $\gamma=\alpha^m$ 为 $p^l$ 次本原单位根, 而 $F=\Q_p(\alpha)=\Q_p(\beta,\gamma)=M_1(\gamma)$. 由 $M_1/\Q_p$ 不分歧可知 $p$ 仍为 $M_1$ 的素元, 而 $\gamma$ 是 $f(x)=\frac{x^{p^l}-1}{x^{p^{l-1}}-1}$ 的根. $\gamma-1$ 是多项式
\[
  g(x)=\frac{(x+1)^{p^l}-1}{(x+1)^{p^{l-1}}-1}
\]
的根. 但是 $g(x)$ 的首项系数为 $1$, 其余系数均被 $p$ 除尽, 而常数项 $g(0)$ 为 $p$. 可知 $g(x)$ 是 $M_1[x]$ 中的 Eisenstein 多项式. 这表明 $F=M_1(\gamma)=M_1(\gamma-1)$ 是 $M_1$ 的 $p^{l-1}(p-1)$ 次完全分歧扩张. 于是 $F/\Q_p$ 是 $p^{l-1}(p-1)s$ 次扩张, $f=s$, $e=p^{l-1}(p-1)$. $F$ 中存在 $t=p^s-1$ 次本原单位根 $\zeta$, $T=\{0,1,\zeta,\cdots,\zeta^{t-1}\}$ 是 $F$ 的赋值环对素理想的陪集完全剩余系, 而 $\gamma-1$ 为 $F$ 的素元. 所以 $F$ 中非零元素 $\delta$ 可惟一表示成
\[
  \delta=\sum_{i=\lambda}^{\infty}c_i(\gamma-1)^i\qquad(c_i\in T,\ c_{\lambda}\neq 0),
\]
其中 $\lambda=v_F(\delta)$.
\end{example}

\begin{Exercises}
\item 设 $E/F$, $K/F$ 为完备离散赋值域的有限扩张, 并且 $\overline{F}$ 是有限域.

(1) 若 $E/F$ 完全分歧, $K/F$ 不分歧, 则 $EK/K$ 完全分歧;

(2) 若 $K/F$ 和 $E/F$ 均完全分歧, 试问 $KE/F$ 是否完全分歧？

(3) 若 $F\subseteq K\subseteq E$, $K/F$ 和 $E/K$ 均完全分歧, 则 $E/F$ 完全分歧.

\item 计算 $\Q_2(\sqrt{3},\sqrt{7})$ 和 $\Q_2(\sqrt{3},\sqrt{2})$ 对于 $\Q_2$ 的分歧指数和剩余类域次数.

\item 设 $F$ 是完备离散赋值域, $\alpha\in F$, $v_F(\alpha)=0$, $\overline{F}$ 是特征为素数 $p$ 的有限域. $n$ 是与 $p$ 互素的正整数, 则 $F(\sqrt[n]{\alpha})$ 是 $F$ 的不分歧扩张.

\item 设 $F$ 是完备离散赋值域, $\alpha\in F$, $v_F(\alpha)\geqslant 1$. 设 $n$ 是与 $v_F(\alpha)$ 互素的正整数. 求证 $F(\sqrt[n]{\alpha})$ 是 $F$ 的 $n$ 次完全分歧扩张.
\end{Exercises}
